% Generated draft: requires mathematical and source review.
\documentclass[11pt]{article}
\usepackage[margin=25mm]{geometry}
\usepackage{fontspec}
\setmainfont{Latin Modern Roman}
\usepackage{amsmath,amssymb,mathtools,unicode-math}
\setmathfont{Latin Modern Math}
\usepackage{enumitem,needspace,xurl,hyperref}
\setlength{\parindent}{0pt}
\setlength{\parskip}{5pt}
\setlength{\emergencystretch}{3em}
\newcommand{\N}{\mathbb{N}}
\newcommand{\Z}{\mathbb{Z}}
\newcommand{\Q}{\mathbb{Q}}
\newcommand{\R}{\mathbb{R}}
\newcommand{\C}{\mathbb{C}}
\newcommand{\F}{\mathbb{F}}
\newcommand{\A}{\mathbb{A}}
\newcommand{\PP}{\mathbb{P}}
\newcommand{\E}{\mathbb{E}}
\newcommand{\eps}{\varepsilon}
\newcommand{\dd}{\,\mathrm{d}}
\begin{document}
\section*{UnsolvedMath ID: 4900010}
\textbf{Source problem code:} catalogue entry
\subsection{Definitions and mathematical statement}
In a standard homogeneous-space setting, let $G$ be a semisimple Lie group, $\Gamma\leq G$ a lattice, and $A\leq G$ a higher-rank abelian diagonalizable subgroup acting on $X=G/\Gamma$ by left multiplication. Here higher rank is intended to exclude a rank-one action; the exact real, arithmetic, or discrete-action setting is not specified in the source. A Borel probability measure $\mu$ is $A$-invariant if $\mu(aB)=\mu(B)$ for every $a\in A$ and Borel set $B$, and $A$-ergodic if every invariant measurable set has measure zero or one.

A homogeneous probability measure is a normalized invariant measure on a closed finite-volume subgroup orbit $Hx\subseteq X$, for a suitable closed subgroup $H\leq G$. The intended meaning of algebraic additionally depends on the allowed subgroup category, which the source does not fix.

Classify the $A$-invariant ergodic probability measures. In particular, specify the intended nondegeneracy hypotheses and prove homogeneity for all measures satisfying them. The phrase ``usual nondegeneracy hypotheses'' in the source does not define a complete quantified conjecture, and higher rank alone is not asserted here to settle every action on every homogeneous space.

The supplied summary attributes substantial positive-entropy rigidity results to Einsiedler--Katok--Lindenstrauss (2006) and Lindenstrauss (2006), while describing the fully general program, especially zero-entropy cases, as open. For an element $a\in A$, positive entropy means $h_\mu(x\mapsto ax)>0$. The summary does not specify the groups, irreducibility assumptions, entropy quantifiers, or other hypotheses of those results, so its broad phrasing is retained as an attribution rather than as a theorem in the entire setting above.
\subsection{Short English statement}
Classify ergodic invariant probability measures for higher-rank diagonal actions on homogeneous spaces. Identify the hypotheses forcing such measures to be homogeneous, while distinguishing known positive-entropy rigidity results from the broader unresolved program.
\subsection{Sources}
\begin{itemize}
\item Wikipedia contributors, \emph{List of unsolved problems in mathematics}, Dynamical systems, current bullet 10 in the supplied extraction: \url{https://en.wikipedia.org/wiki/List_of_unsolved_problems_in_mathematics#Dynamical_systems}; accessed 2026-07-29, HTML extraction. The record attributes the program to Grigory Margulis. Wikipedia is not an authoritative resolution source.
\item Supplied research-summary attributions: Einsiedler--Katok--Lindenstrauss (2006) and Lindenstrauss (2006), positive-entropy measure-rigidity results. Titles, theorem numbers, and full hypotheses are not supplied.
\item Local registry \texttt{data/unsolvedmath/problems.json}, permanent UnsolvedMath ID $4900010$, source code catalogue entry, \emph{Margulis measure-classification conjecture}. Dataset attribution: Ulam AI and the UnsolvedMath Contributors.
\item The supplied Wikipedia rights notice states CC BY-SA 4.0 and requires attribution and share-alike compatibility to be addressed in release review.
\end{itemize}
\end{document}
